\section{The hole relation and its tensor realization}\label{sec:geometry}

Three pairwise missing edges would give an independent set of size
three. We choose a relation that excludes such a triangle by parity.
Later we choose its linear maps so that four missing edges occur often
enough between two units.

\subsection{A triangle-free relation}

Let \(\cX\) and \(\cV\) be finite-dimensional vector spaces over
\(\F\). Fix a linear functional \(a\in\cX^*\) and a symmetric
bilinear form \(T\) on \(\cX\), and write
\[
 r(x)=a+T(x,\cdot)\in\cX^*.
\]
Let \(\Omega\) be a finite set. For each \(i\in\Omega\), suppose
we have an injective linear map \(U_i:\cX\to\cV\) and a linear
functional \(u_i\in\cV^*\).

% lax begin Lax793645.Holes
\begin{definition}\label{def:hole}
Two elements \(i,j\in\Omega\) have a \emph{hole} between them if
there exist \(\lambda_i,\lambda_j\in\cX\) such that
\begin{equation}\label{eq:hole-equations}
 \begin{gathered}
 U_i\lambda_i=U_j\lambda_j,\qquad
 u_jU_i=r(\lambda_i),\qquad u_iU_j=r(\lambda_j),\\
 a(\lambda_i)+a(\lambda_j)=1.
 \end{gathered}
\end{equation}
The middle two identities are identities of linear functionals on
the whole space \(\cX\).
\end{definition}
% lax end

Thus a hole has a shared vector in \(\cV\), one witness for it in
each copy of \(\cX\), and opposite values of \(a\) on the witnesses.
The functional identities make these requirements incompatible around
a triangle.

% lax begin Lax793645.TriangleFreeHoles
\begin{lemma}\label{lem:hole-triangle-free}
The hole relation is symmetric, has no loops, and is triangle-free.
\end{lemma}
% lax end
% lax begin Lax793645Proofs.Manuscript.TriangleFreeHoles.triangle_free
\begin{proof}
Symmetry follows by interchanging the endpoints.
A loop would have \(U_i\lambda_i=U_i\lambda_j\), hence
\(\lambda_i=\lambda_j\) by injectivity, contradicting the last
equation of \Cref{eq:hole-equations}.

Suppose that three distinct elements have holes on all three pairs.
For a directed incidence \(ij\), let \(x_{ij}\) be its witness at
\(i\). For distinct \(i,j,k\), evaluating the functional identity for
\(ij\) on \(x_{ik}\) gives
\[
 u_j(U_ix_{ik})=a(x_{ik})+T(x_{ij},x_{ik}).
\]
Sum over the six ordered choices of \((i,j,k)\).
For fixed \(j\), the two arguments on the left are equal by the
sharing equation for the remaining edge \(ik\), so they cancel.
For fixed \(i\), the two bilinear terms cancel by symmetry of \(T\).
The remaining sum is
\[
 \sum_{\{i,j\}}\bigl(a(x_{ij})+a(x_{ji})\bigr)=1+1+1=1,
\]
a contradiction.
\end{proof}
% lax end

Given any list \(o_1,\ldots,o_m\in\Omega\), form a graph \(G\)
on its \emph{positions}: two distinct positions are adjacent when
their elements have no hole. This gives a finite simple graph even
when elements repeat. Equal elements are adjacent because holes have
no loops. An independent triple of positions would therefore give
three distinct elements forming a hole triangle. Hence
\begin{equation}\label{eq:sample-independence}
 \alpha(G)\le2,\qquad \chi(G)\ge\lceil m/2\rceil.
\end{equation}

It remains to choose the finite set and its linear data so that a
random list also has no large connected matching. Two disjoint edges
fail to touch precisely when all four cross pairs are holes. The
distribution theorem in \Cref{sec:distributions} will force these
four-hole conflicts for the model constructed below.


\subsection{Points and the odd cycle}

Fix a large odd integer \(g\). Tags are the elements of \(\mathbb Z/g\mathbb Z\).
There is one component
\[
 e_l=\{l,l+1\}
\]
for each tag \(l\). We choose \(g\) after the initial pin budget, as
specified in Appendix~\ref{sec:parameters}. It does not grow with the graph.

The base coordinates have one constant coordinate and blocks
\[
 O_{l,t}\quad(l,t\in\mathbb Z/g\mathbb Z),\qquad S,\quad \#,\quad Z.
\]
Each block has \(n\) coordinates. At tag \(l\), only the ordinary blocks
\(O_{l,t}\) are allowed. The blocks \(S,\#,Z\) are shared by all tags.
We write \(F\) for the base vector space, including its constant coordinate.
A base point is \(\bar z=(1,z)\in F\). A point allowed at tag \(l\)
has zero coordinates in every \(O_{l',t}\) with \(l'\ne l\).

Labels are vectors \(s\in\F^b\). Set
\[
 \mathsf S=\F^{b+1},\qquad p_s=(1,s),\qquad
 B=\mathsf S\otimes F,\qquad v(s,z)=p_s\otimes\bar z.
\]
Later, \(S_0\) will denote a protected subspace of \(\mathsf S\).
The point matrix is \(v(s,z)v(s,z)\trans\).
Let \(W_l\) be the span of point matrices allowed at tag \(l\).
Let \(W_*\) be the corresponding span when all ordinary coordinates are
zero. Then \(W_l\cap W_{l'}=W_*\) for distinct tags. To see this, use
the ordinary affine point-matrix basis: its entries involve only coordinates
allowed at that tag. The entries common to two tags are exactly the shared
ones.

Define the space of profiles by
\[
 \cX=\left\{x:\ x_{e_l}=w_l+w_{l+1},\quad w_l\in W_l\right\}.
\]
Thus a point at tag \(l\) gives a profile, called an \emph{atom}, with its
point matrix on the two incident components and zero elsewhere.
We denote it by \(A(l,s,z)\). The kernel of \((w_l)_l\mapsto x\)
consists of constant tuples with value in \(W_*\): equality of consecutive
terms propagates around the cycle.

\subsection{Testers and the scalar form}

In each ordinary block and in \(S\), reserve \(2r_0\) coordinates.
Pair them and sum their products. These quadratic expressions define linear
functionals on point matrices, denoted by \(\eta_{l,t}\) and \(\eta_S\).
Write
\[
 \eta_O=\sum_{l,t}\eta_{l,t},\qquad \eta=\eta_O+\eta_S.
\]
All tester functionals ignore \(\#\) and \(Z\). Their matrices use the
constant selector coordinate. The polar rank of a sum of \(a\) disjoint
testers is \(2ar_0\).

Delete tag \(t\) from the cycle. The remaining path has the perfect matching
\[
 M_t=\{e_{t+1},e_{t+3},\ldots,e_{t+g-2}\}.
\]
Define linear functionals directly on the component matrices:
\begin{equation}\label{eq:cycle-functionals}
 a_t(x)=\sum_l\eta_{l,t}(x_{e_l}),\qquad
 a(x)=\sum_ta_t(x),\qquad
 b_t(x)=\sum_{e\in M_t}\eta_S(x_e).
\end{equation}
If \(x_{e_l}=w_l+w_{l+1}\), the private ordinary blocks and the matching
identity give
\[
 a_t(x)=\sum_l\eta_{l,t}(w_l),\qquad
 b_t(x)=\sum_{l\ne t}\eta_S(w_l).
\]
These formulas explain the definitions; no choice of representation is
needed to define the functionals.

The first part of the scalar form is
\[
 T^0(x,y)=a(x)a(y)+\sum_t\bigl(a_t(x)b_t(y)+b_t(x)a_t(y)\bigr).
\]
In particular \(T^0(x,x)=a(x)\).
For later use, put \(z_c(x)=\sum_{e\in M_c}x_e\).
Choose matrices \(L_j^c,R_j^c\), with \(1\le j\le J\), and define
\begin{equation}\label{eq:reduced-mixer}
 B^1(x,y)=\sum_{c,j}\pair{x_{e_c}}{L_j^c z_c(y)(R_j^c)\trans},
 \qquad T^1(x,y)=B^1(x,y)+B^1(y,x).
\end{equation}
Here \(\pair{A}{B}=\sum_{i,j}A_{ij}B_{ij}\). Set \(T=T^0+T^1\).
This form is symmetric and satisfies \(T(x,x)=a(x)\).
The matrices will be chosen in \Cref{sec:moments}. There are \(Jg\)
pairs of them.

The matching has a useful effect on atoms:
\begin{equation}\label{eq:matching-atom}
 z_c(A(l,s,z))=
 \begin{cases}v(s,z)v(s,z)\trans,&c\ne l,\\0,&c=l.\end{cases}
\end{equation}
Every tag other than \(c\) occurs exactly once in \(M_c\).
This is why a small mixer family suffices.

\subsection{Counting ordinary tester bits}

Put \(o_l(x)=\sum_t\eta_{l,t}(x_{e_l})\). Privacy of the ordinary
blocks gives
\[
 o_l(x)+o_{l+1}(x)=\eta_O(x_{e_l}),\qquad
 \sum_l o_l(x)=a(x).
\]
We need only the following support bound.
% lax begin Lax793645.OrdinaryTesterSupport
\begin{lemma}\label{lem:ordinary-support}
For every profile \(x\), \(|\supp o(x)|\le\sum_e\rank x_e\).
\end{lemma}
% lax end
% lax begin Lax793645Proofs.Manuscript.OrdinaryTesterSupport.ordinaryBit_support
\begin{proof}
If \(o_l(x)=1\), an ordinary entry owned by \(l\) is nonzero.
Both incident components contain that entry. Each such tag contributes
two incidences with nonzero components, and each component receives at
most two. Thus the support size is at most the number of nonzero
components, which is at most their total rank.
\end{proof}
% lax end

\subsection{One quadratic frame per component}

Choose a matrix \(E\) on \(B\) whose value on each point is \(\eta\).
It consists of the tester matrices and a correction \(E^\#\) supported
on \(\#\). The latter has zero value on a single point; its off-point
values will ensure binary mixing. Define \(q_E(v)=v\trans Ev\) for
\emph{every} \(v\in B\).

Let \(V=\F^{N/2}\oplus\F^{N/2}\), with
\[
 q_0(x,y)=x\cdot y,
 \qquad \beta_0(v,w)=q_0(v+w)+q_0(v)+q_0(w).
\]
At each component choose an injective linear map on
\(B\oplus\F^h\oplus\F^h\). Its restrictions are denoted by \(P,X,Y\).
They satisfy
\begin{equation}\label{eq:one-frame}
 q_0(Pv)=q_E(v),\qquad \beta_0(X_j,Pv)=0
 \quad(v\in B,\ 1\le j\le h).
\end{equation}
There are no other prescribed quadratic values or pairings on the channel
columns. The full list of primal and channel columns is independent.
The uniform law on these frames exists once \(N\) is sufficiently large;
\Cref{sec:frame-laws} proves this and the counting facts we need.

A raw vertex \(o\) is an independent choice of such a frame at every
component. Its law is \(\mu\), on a finite set \(\Omega_n\).
Let \(\cV\) be the product of the ambient symmetric-matrix spaces. Define
\[
 (U_ox)_e=P_ex_eP_e\trans.
\]
The quadratic form
\[
 v\longmapsto\sum_{j=1}^h\beta_0(X_{e,j},v)\beta_0(Y_{e,j},v)
\]
defines a linear functional on symmetric matrices. Sum it over the
components to obtain \(u_o\in\cV^*\).
A left inverse to each \(P_e\) proves that \(U_o\) is injective.
Equation \eqref{eq:one-frame} gives \(u_oU_o=0\).
We use these maps in Definition~\ref{def:hole}.

The ambient functional defined by \(q_0\) satisfies
\[
 q_0[P_ex_eP_e\trans]=\eta(x_e)\quad(x_e\in W_l+W_{l+1}).
\]
Square brackets denote evaluation on a symmetric matrix. Therefore
\(U_ox=U_{o'}y\) implies equality of the component tester values:
\begin{equation}\label{eq:common-image-tester}
 U_ox=U_{o'}y\quad\Longrightarrow\quad\eta(x_e)=\eta(y_e)\quad\text{for every }e.
\end{equation}

All coefficient dimensions are \(O(n)\), with fixed implied constants.
Counting the entries of the component frames gives
\begin{equation}\label{eq:raw-space-size}
 \log_2|\Omega_n|=O(N^2).
\end{equation}
For the graph transfer we use
\[
 D=8000,\qquad M=2^{1000},\qquad m=2^{2000N},\qquad
 \varepsilon_N=2^{-200N}.
\]
The remaining choices, including \(h=1000r_0\) and \(M_0\), are made
in the order recorded in Appendix~\ref{sec:parameters}.
